\section{Physical Experimental Results}
\label{sec:experimental_results}

\subsection{Experimental Data Scope}

The physical experiments were performed on the STM32F411--ESP-NOW--ROS\,2
platform using overhead AprilTag localization fused with wheel odometry.
Figure~\ref{fig:exp_setup} shows the camera view and the marked robot in the
test area. All reported errors were recomputed from the camera pose after
rigid alignment to the local trajectory frame; no temporal smoothing was
applied during post-processing.

\begin{figure}[!t]
  \centering
  \includegraphics[width=\columnwidth]
    {\ExpDir/platform_overhead_camera.jpeg}
  \caption{Overhead-camera view of the physical test area. The AprilTag
  attached to the robot provides the absolute pose measurement used by the
  EKF.}
  \label{fig:exp_setup}
\end{figure}

The currently available data do not form a balanced final validation set.
The circular nominal runs contain five laps per controller; the selected
square runs contain three BSMC laps, two Backstepping laps, and one SMC
tuning lap; and the figure-eight comparison uses the best available pilot
run for each controller. Consequently, the results below are presented as
an experimental pilot comparison. In particular, the SMC square run is
tuning data and is not relabeled as validation data.

\subsection{Nominal Trajectory Tracking}

Figure~\ref{fig:exp_nominal_tracking} compares the measured trajectories for
the three reference families. On the circle, all controllers remain close to
the reference, with nominal Backstepping giving the lowest selected position
RMSE. On the figure-eight, BSMC and SMC give similar position RMSE, whereas
BSMC gives the lowest heading RMSE. For the sharp-corner square, BSMC and
Backstepping are close, while the selected pure-SMC pilot exhibits larger
outward corner transients.

\begin{figure*}[!t]
  \centering
  \subfloat[Circle: five-lap selected runs.]{%
    \includegraphics[width=0.31\textwidth]
      {\ExpDir/circle_nominal/circle_three_controller_trajectory.pdf}%
    \label{fig:exp_circle_tracking}}
  \hfil
  \subfloat[Figure-eight: selected pilot runs.]{%
    \includegraphics[width=0.31\textwidth]
      {\ExpDir/eight_nominal/eight_three_controller_trajectory.pdf}%
    \label{fig:exp_eight_tracking}}
  \hfil
  \subfloat[Square: selected runs with unequal lap counts.]{%
    \includegraphics[width=0.31\textwidth]
      {\ExpDir/square_nominal/square_three_controller_trajectory.pdf}%
    \label{fig:exp_square_tracking}}
  \caption{Measured physical trajectory tracking for Backstepping, BSMC,
  and SMC. Each trajectory is expressed in its rigidly aligned local frame.
  The square SMC trace is a tuning pilot rather than a final validation run.}
  \label{fig:exp_nominal_tracking}
\end{figure*}

\begin{table}[!t]
  \centering
  \caption{Selected Physical Tracking Runs}
  \label{tab:exp_tracking}
  \scriptsize
  \setlength{\tabcolsep}{3.2pt}
  \renewcommand{\arraystretch}{1.08}
  \begin{tabular}{llcc}
    \toprule
    \textbf{Path} & \textbf{Controller}
      & \textbf{Position RMSE} & \textbf{Heading RMSE}\\
      & & (cm) & (deg)\\
    \midrule
    \multirow{3}{*}{Circle}
      & Backstepping & \textbf{2.12} & \textbf{1.28}\\
      & BSMC         & 2.27 & 1.52\\
      & SMC          & 2.53 & 2.00\\
    \midrule
    \multirow{3}{*}{Eight}
      & Backstepping & 3.16 & 4.69\\
      & BSMC         & 2.68 & \textbf{3.33}\\
      & SMC          & \textbf{2.66} & 5.32\\
    \midrule
    \multirow{3}{*}{Square}
      & Backstepping & 3.22 & \textbf{22.26}\\
      & BSMC         & \textbf{3.04} & 22.59\\
      & SMC$^\dagger$& 8.95 & 24.12\\
    \bottomrule
  \end{tabular}

  \vspace{1mm}
  \raggedright
  \footnotesize $^\dagger$Tuning pilot; not a locked validation run.
\end{table}

The large overall square heading RMSE values arise primarily from the
instantaneous $90^\circ$ reference-heading changes at the vertices. To
separate corner transients from established straight-edge behavior,
Table~\ref{tab:exp_square_straights} reports metrics only on the detected
straight portions. BSMC gives the lowest straight lateral RMSE, while
Backstepping gives the lowest straight heading RMSE. The SMC pilot requires
substantially more angular correction and does not return to the new edge as
quickly after each corner.

\begin{table}[!t]
  \centering
  \caption{Straight-Edge Metrics for the Selected Square Runs}
  \label{tab:exp_square_straights}
  \scriptsize
  \setlength{\tabcolsep}{2.4pt}
  \renewcommand{\arraystretch}{1.08}
  \begin{tabular}{lcccc}
    \toprule
    \textbf{Controller}
      & \shortstack{\textbf{Lateral}\\\textbf{RMSE}}
      & \shortstack{\textbf{Lateral}\\\textbf{bias}}
      & \shortstack{\textbf{Maximum}\\$|e_y|$}
      & \shortstack{\textbf{Heading}\\\textbf{RMSE}}\\
      & (cm) & (cm) & (cm) & (deg)\\
    \midrule
    Backstepping & 3.44 & 3.28 & \textbf{5.35} & \textbf{2.11}\\
    BSMC         & \textbf{3.26} & \textbf{3.07} & 5.53 & 2.24\\
    SMC$^\dagger$& 8.94 & 6.79 & 19.19 & 7.90\\
    \bottomrule
  \end{tabular}
\end{table}

\begin{figure*}[!t]
  \centering
  \subfloat[Circle error histories.]{%
    \includegraphics[width=0.31\textwidth]
      {\ExpDir/circle_nominal/circle_three_controller_errors.pdf}}
  \hfil
  \subfloat[Figure-eight error histories.]{%
    \includegraphics[width=0.31\textwidth]
      {\ExpDir/eight_nominal/eight_three_controller_errors.pdf}}
  \hfil
  \subfloat[Square: common first-lap time window.]{%
    \includegraphics[width=0.31\textwidth]
      {\ExpDir/square_nominal/square_three_controller_errors.pdf}}
  \caption{Measured position and heading error histories for the selected
  nominal physical runs. The periodic square heading peaks coincide with the
  commanded $90^\circ$ corner transitions.}
  \label{fig:exp_nominal_errors}
\end{figure*}

\subsection{Manual Disturbance and Relocalization Pilots}

The manual disturbance protocol was applied only to the circular trajectory.
The robot was paused, physically displaced, and then released while the
controller continued from the same reference phase. Two displacement
orientations were recorded: nominal yaw ($0^\circ$) and an approximately
$90^\circ$ yaw change. These trials were designed as relocalization and safety
pilots, not as balanced repeat statistics.

\begin{figure*}[!t]
  \centering
  \subfloat[Yaw-$0^\circ$ trajectories aligned to a common intervention
  phase.]{%
    \includegraphics[width=0.31\textwidth]
      {\ExpDir/disturbance_yaw0/circle_three_controller_disturbance_trajectory.pdf}}
  \hfil
  \subfloat[Yaw-$0^\circ$ recovery histories.]{%
    \includegraphics[width=0.31\textwidth]
      {\ExpDir/disturbance_yaw0/circle_three_controller_disturbance_recovery.pdf}}
  \hfil
  \subfloat[Yaw-$90^\circ$ recovery histories.]{%
    \includegraphics[width=0.31\textwidth]
      {\ExpDir/disturbance_yaw90/circle_three_controller_disturbance_recovery.pdf}}
  \caption{Physical manual-disturbance pilots on the circular path. For panel
  (a), each measured trajectory is phase-normalized by a smoothly varying
  rotation and translation between the common run-start and BSMC intervention
  landmarks; the transform is fixed after intervention and no measurement
  smoothing is applied. A missing recovery value denotes failure to satisfy
  the fixed recovery criterion within the recorded interval. The selected
  yaw-$90^\circ$ SMC trial was safety-censored.}
  \label{fig:exp_disturbance}
\end{figure*}

\begin{table}[!t]
  \centering
  \caption{Selected Circle Disturbance Pilots}
  \label{tab:exp_disturbance}
  \scriptsize
  \setlength{\tabcolsep}{2.4pt}
  \renewcommand{\arraystretch}{1.08}
  \begin{tabular}{llccc}
    \toprule
    \textbf{Yaw} & \textbf{Controller}
      & \textbf{Peak increment} & $\mathbf{IAE_{10}}$
      & $\mathbf{T_{\mathrm{rec}}}$\\
      & & (cm) & (m\,s) & (s)\\
    \midrule
    \multirow{3}{*}{$0^\circ$}
      & Backstepping & 30.17 & 1.339 & 6.88\\
      & BSMC         & 28.37 & \textbf{1.279} & 7.24\\
      & SMC          & \textbf{19.24} & 1.366 & NR\\
    \midrule
    \multirow{3}{*}{$90^\circ$}
      & Backstepping & \textbf{52.66} & \textbf{3.451} & NR\\
      & BSMC         & 54.19 & 3.778 & NR\\
      & SMC$^\ddagger$ & 74.40 & 3.516 & NR\\
    \bottomrule
  \end{tabular}

  \vspace{1mm}
  \raggedright
  \footnotesize NR: recovery criterion not reached. 
  $^\ddagger$Safety-censored run.
\end{table}

The yaw-$0^\circ$ pilot shows successful recovery for Backstepping and BSMC,
whereas the selected SMC run does not meet the recovery criterion. None of
the selected yaw-$90^\circ$ trials meets the same criterion. The result
demonstrates that large pose relocalization is a substantially harder
closed-loop condition than nominal tracking. It does not establish a
statistical ordering because only selected pilot runs are available and the
applied manual displacements are not identical.

\subsection{Implementation Observations and Validity}

The experiments revealed two implementation effects that are not visible in
the ideal kinematic simulation. First, absolute camera measurements must be
allowed to re-enter the EKF after a sustained large innovation; otherwise the
filter can continue from an encoder-only phantom pose after the robot is
physically displaced. Second, square-corner performance depends strongly on
the inner wheel-speed loop. For the same nominal corner command, a stalled
inner wheel increases the effective turning radius and leaves less straight
distance for lateral-error recovery. These effects motivate reporting
camera-derived errors and actuator telemetry in addition to controller
commands.

The nominal circle data are the most balanced set presently available. The
figure-eight and square comparisons remain selected pilots, and the
disturbance data are manual safety trials. Final publication claims requiring
mean and standard deviation must therefore be based on a future locked-gain,
equal-repeat validation campaign rather than on the selected-run values in
this section.

\FloatBarrier
